% Prostudováno: ucebnice/02_Funkce_a_rovnice/Linearni_funkce_a_rovnice/Linerání rovnice.tex,
% Soustavy lineárních rovnic, pododdíl Sčítací metoda (ř. 1123–1218), a kontrola společné dvojice.
% Realisticky.cz: externi_zdroje/realisticky/02_Funkce_a_rovnice/03_Nektere_rovnice_a_nerovnice_prevoditelne_na_linearni/020310_Soustavy_dvou_linearnich_rovnic_o_dvou_neznamych_I/020310_Lekce.pdf
% a 020311_Soustavy_dvou_linearnich_rovnic_o_dvou_neznamych_II/020311_Lekce.pdf; využito srovnání metod a významu eliminace. Úlohy vlastní.
\documentclass[11pt,a4paper]{article}
% Spolecny styl pro ucitelske pripravy MAT (Work/MAT)
% Verze 1.1, 2026-10-23
\NeedsTeXFormat{LaTeX2e}
\ProvidesPackage{priprava}[2026/10/23 v1.1 Sablona ucitelskych priprav MAT]

% Zakladni typografie a cestina
\RequirePackage[T1]{fontenc}
\RequirePackage[utf8]{inputenc}
\RequirePackage[czech]{babel}
\RequirePackage{lmodern}
\RequirePackage{microtype}
\RequirePackage{amsmath,amssymb}
\RequirePackage{xcolor}
\RequirePackage{enumitem}
\RequirePackage[most]{tcolorbox}
\RequirePackage{tabularx,booktabs,array}
\RequirePackage{fancyhdr}
\RequirePackage{needspace}
\RequirePackage{geometry}
\RequirePackage[hidelinks]{hyperref}

\geometry{a4paper,top=17mm,bottom=18mm,left=23mm,right=23mm}
\definecolor{notegray}{gray}{0.39}
\setlength{\parindent}{0pt}
\setlength{\parskip}{2.5pt plus 1pt minus 0.5pt}
\setlist[itemize]{topsep=2pt,itemsep=1pt,parsep=0pt,leftmargin=1.45em}
\setlist[enumerate]{topsep=2pt,itemsep=2pt,parsep=0pt,leftmargin=1.65em}
\setcounter{secnumdepth}{0}
\renewcommand{\arraystretch}{1.12}

% Zahlavi: vlevo rocnik | globalni tema, vpravo cislo: nazev hodiny
% Zapati zustava zcela prazdne.
\pagestyle{fancy}
\fancyhf{}
\renewcommand{\headrulewidth}{0.35pt}
\renewcommand{\footrulewidth}{0pt}
\setlength{\headheight}{14pt}
\newcommand{\nastavPripravu}[4]{%
  \fancyhead[L]{\small #1\enspace\textbar\enspace #2}%
  \fancyhead[R]{\small #3: #4}%
}

% Minimalisticke cernobile boxy. Kazda definice zvlast.
\tcbset{commonlesson/.style={enhanced,breakable=false,colback=white,colframe=black,
  colbacktitle=white,coltitle=black,boxrule=0.45pt,
  arc=0.8mm,outer arc=0.8mm,left=3mm,right=3mm,top=1.6mm,bottom=1.6mm,
  toptitle=1mm,bottomtitle=0.7mm,fonttitle=\bfseries,
  attach boxed title to top left={xshift=2.5mm,yshift=-1.15mm},
  boxed title style={colback=white,colframe=white,boxrule=0pt,left=1mm,right=1mm}}}
\newtcolorbox{zadaniBox}[1][]{commonlesson,title={Zadání},#1}
\newtcolorbox{prikladBox}[1][]{commonlesson,title={Příklad},#1}
\newtcolorbox{definiceBox}[1][]{commonlesson,title={Definice},#1}

% Spolecny box pro zadani a jeho reseni, bez mezery mezi castmi.
% Pouziti:
%   \begin{prikladReseniBox}
%     Zadani ulohy.
%     \tcblower
%     \begin{align*} ... \end{align*}
%   \end{prikladReseniBox}
% Spodni cast ma jemny sedy podklad a delici linku; lze v ni sazet
% vice radku matematiky. Obe casti maji jen jeden spolecny ramecek.
\tcbset{solutionlesson/.style={
  commonlesson,
  enhanced jigsaw,
  breakable,
  colbacklower=black!3,
  segmentation style={solid,draw=black!25,line width=0.35pt},
  fontlower=\small\color{notegray},
  before lower={\noindent\textbf{Řešení:}\enspace},
  leftlower=3mm,rightlower=3mm
}}
\newtcolorbox{prikladReseniBox}[1][]{solutionlesson,title={Příklad},#1}
\newtcolorbox{zadaniReseniBox}[1][]{solutionlesson,title={Zadání},#1}

% Pomucky pro autora pripravy: sedive poznamky; \reseni pro starsi pripravy
\newcommand{\poznamka}[1]{\par\begingroup\small\color{notegray}\itshape
  \textnormal{\textbf{Poznámka pro učitele.}}\enspace #1\par\endgroup}
\newcommand{\reseni}[1]{\par\begingroup\small\color{notegray}
  \textbf{Očekávané řešení:}\enspace #1\par\endgroup}
\newcommand{\blok}[1]{\needspace{5\baselineskip}\subsection*{#1}}
\newcommand{\vyrok}[1]{\uv{#1}}

\nastavPripravu{1. ročník}{Soustavy rovnic}{78}{Sčítací metoda}
\begin{document}
\blok{Odstranění jedné neznámé součtem}
V minulé hodině jsme jednu neznámou vyjádřili a dosadili. Někdy jsou koeficienty u jedné neznámé opačné, takže při sečtení rovnic tato neznámá zmizí. Sčítáme celé levé strany i celé pravé strany; upravujeme-li nejprve rovnici násobením, násobíme každou její část.
\begin{definiceBox}
\textbf{Sčítací metoda} spočívá v ekvivalentním vynásobení jedné či obou rovnic nenulovým číslem tak, aby koeficienty u zvolené neznámé byly opačné. Součet rovnic pak dá rovnici s jednou neznámou; druhou hodnotu dopočteme z některé původní rovnice a ověříme obě.
\end{definiceBox}
\begin{prikladReseniBox}
Řešte soustavu $2x+y=7$, $3x-y=8$ sčítací metodou a ověřte výsledek.
\tcblower
Koeficienty u $y$ jsou $1$ a $-1$. Sečtením dostaneme $5x=15$, tedy $x=3$. Z první rovnice $2\cdot3+y=7$, takže $y=1$. Zkouška: $2\cdot3+1=7$ a $3\cdot3-1=8$. Proto $K=\{(3;1)\}$.
\end{prikladReseniBox}
\begin{prikladReseniBox}
Řešte $3x+2y=12$, $2x-3y=-5$. Připravte opačné koeficienty u $y$.
\tcblower
První rovnici násobíme $3$ a druhou $2$:
\[
9x+6y=36,\qquad 4x-6y=-10.
\]
Součet dává $13x=26$, tedy $x=2$. Z původní první rovnice $6+2y=12$, takže $y=3$. Zkouška: $3\cdot2+2\cdot3=12$ a $2\cdot2-3\cdot3=4-9=-5$; $K=\{(2;3)\}$.
\end{prikladReseniBox}
\poznamka{Sečtená rovnice sama o sobě nenahrazuje celou soustavu: určí jednu hodnotu a druhou musíme dopočítat i ověřit. Při násobení rovnice číslem $-1$ se mění znaménka všech členů včetně pravé strany.}
\blok{Rozhodnutí mezi metodami}
Dosazování je obvykle rychlé při koeficientu $1$ či $-1$ u některé neznámé. Sčítání je výhodné, když jsou koeficienty již opačné nebo se snadno stanou opačnými. Metody mají dát stejnou množinu řešení; volba je otázkou úsporného postupu, nikoli jiné matematiky.
\begin{prikladReseniBox}
Vyřešte $x+2y=9$, $2x-2y=0$ sčítáním a pak v jedné větě popište dosazovací kontrolu.
\tcblower
Součet dává $3x=9$, tedy $x=3$. Z druhé rovnice $6-2y=0$, takže $y=3$. Zkouška: $3+6=9$, $6-6=0$. Dosazovací kontrola: z druhé rovnice $x=y$, vložení do první dává $3y=9$.
\end{prikladReseniBox}
\newpage
\blok{Samostatné procvičení}
\begin{zadaniReseniBox}
\begin{enumerate}[label=\alph*)]
\item Řešte $x+y=6$, $x-y=2$.
\item Řešte $2x+3y=12$, $2x-y=4$.
\item Řešte $3x-y=8$, $x+2y=5$; zvolte násobení, které odstraní $y$.
\end{enumerate}
\tcblower
\begin{enumerate}[label=\alph*)]
\item Součet: $2x=8$, $x=4$, $y=2$. Kontrola $4+2=6$, $4-2=2$.
\item Odečteme druhou rovnici od první: $4y=8$, $y=2$. Pak $2x-2=4$, $x=3$. Kontrola $6+6=12$, $6-2=4$.
\item Druhou rovnici ponecháme, první vynásobíme $2$: $6x-2y=16$ a $x+2y=5$. Součet $7x=21$, $x=3$, $y=1$. Kontrola $9-1=8$, $3+2=5$.
\end{enumerate}
\end{zadaniReseniBox}
\blok{Volitelné rozšíření a rezerva}
\begin{zadaniReseniBox}
Porovnejte soustavy $2x-3y=1$, $4x-6y=2$ a $2x-3y=1$, $4x-6y=3$. Co ukáže eliminace?
\tcblower
Po odečtení dvojnásobku první rovnice od druhé dostaneme v prvním případě $0=0$: druhá podmínka je totožná s první. Nekonečně mnoho řešení lze zapsat $K=\{((1+3t)/2;t):t\in\mathbb R\}$. Ve druhém případě dostaneme $0=1$, tedy $K=\varnothing$. Výsledek eliminace je nutno interpretovat, ne dělit nulovým koeficientem.
\end{zadaniReseniBox}
\blok{Domácí práce}
\begin{zadaniReseniBox}
\begin{enumerate}[label=\alph*)]
\item Řešte $5x+2y=16$, $3x-2y=0$.
\item Řešte $2x+y=9$, $4x-2y=6$.
\item Vysvětlete, proč z $2x+y=7$ a $3x-y=8$ lze odvodit $5x=15$, avšak pro kontrolu dvojice nestačí ověřit jen poslední rovnici.
\end{enumerate}
\tcblower
\begin{enumerate}[label=\alph*)]
\item Sečtením $8x=16$, tedy $x=2$, $y=3$. Kontrola $10+6=16$, $6-6=0$.
\item První rovnici vynásobíme $2$: $4x+2y=18$. Součet s druhou dává $8x=24$, $x=3$, $y=3$. Kontrola $6+3=9$, $12-6=6$.
\item Součet původních rovnic je nutnou podmínkou, ale $5x=15$ neobsahuje informaci o $y$. Např. $(3;0)$ ji splňuje, avšak první původní rovnici ne ($6\ne7$). Proto druhou hodnotu dopočteme a ověříme obě rovnice.
\end{enumerate}
\end{zadaniReseniBox}
\poznamka{Jádro: princip 8 min, dva vzory 15 min, samostatná práce a kontrola 19 min, organizace 3 min. Třetí vzor a rezerva podle tempa; domácí práce 15–20 min. Příští hodina spojí počet řešení s polohou přímek.}
\end{document}
